\begin{center}
  {\small AVENGERS INITIATIVE -- NEW RECRUIT ORIENTATION}\par
  \vspace{0.3em}
  {\Large\bfseries LESSON PLAN: GREAT POWER, GREATER RESPONSIBILITY}\par
  \vspace{0.2em}
  {\normalsize \instructorlabel\ Peter Parker / Spider-Man}
\end{center}

\dossierbox{\identificationtitle}{\identificationfields}

\dossierbox{\objectivetitle}{%
\textbf{General:} Develop responsible decision-making by considering the needs of others, checking what is known, and anticipating the consequences of action during a fictional emergency.\par
\textbf{Specific objectives:}
\begin{dossieritems}
  \item Identify people affected by the incident and describe their immediate safety needs;
  \item Distinguish verified information from uncertain reports and identify what should be checked next;
  \item Propose and justify a coordinated response, including one likely consequence, limitation, or reason to seek additional help.
\end{dossieritems}
}

\dossierbox{\contenttitle}{%
\begin{dossieritems}
  \item Responsibility, care for others, and the choices that come with power;
  \item Community needs, reliable information, and uncertainty during an emergency;
  \item Team communication, safe coordination, and the consequences of a response.
\end{dossieritems}
}

\dossierbox{\methodologytitle}{%
Peter Parker leads a guided, case-based role-play for new Avengers recruits. Teams receive the same fictional subway incident but begin with incomplete and conflicting reports. They examine who may be affected, assign complementary responsibilities, and plan a response they can explain. A new piece of information is introduced during the activity so teams can reconsider their assumptions. The instructor evaluates the care, evidence, communication, and safety behind each decision rather than expecting one scripted answer.
}

\dossierbox{\procedurestitle}{%
\textbf{Opening and briefing (5 minutes):} Spider-Man introduces a citywide blackout and a malfunctioning support drone near a crowded subway entrance. The first reports are incomplete, and passengers may need assistance. Recruits review the mission card, the 30-minute schedule, and the rule to prioritize people's safety and request qualified help when needed.\par
\textbf{Development (20 minutes):} Teams assign roles such as coordinator, passenger advocate, information checker, and safety lead. They map who may be affected and separate confirmed details from open questions (5 minutes). Each team then plans a safe sequence of actions and explains how members will communicate and divide responsibilities (10 minutes). The instructor shares a new report; teams decide whether it changes their plan and record one reason for their decision (5 minutes).\par
\textbf{Integration and conclusion (5 minutes):} Teams briefly share one action, the evidence supporting it, and a remaining risk. Spider-Man connects their choices to responsibility and care for the community. Each recruit completes an exit reflection naming one responsible action and one question that should be verified next.
}

\dossierbox{\resourcestitle}{%
\begin{dossieritems}
  \item Projector or display for the fictional mission briefing;
  \item Printed mission cards and team-role cards;
  \item Whiteboard and markers for mapping affected people and verified information;
  \item Response-planning sheet and individual exit-reflection card for each recruit.
\end{dossieritems}
}

\dossierbox{\assessmenttitle}{%
Assessment is formative and based on team discussion, the response-planning sheet, and each recruit's exit reflection. Evidence of learning includes the recruit's ability to:
\begin{dossieritems}
  \item Describe affected people's needs and prioritize safety;
  \item Separate confirmed information from reports that still need verification;
  \item Explain a coordinated response and communicate responsibilities;
  \item Identify a likely consequence, limitation, or reason to seek qualified help.
\end{dossieritems}
}

\dossierbox{\bibliographytitle}{%
\begin{dossieritems}
  \item Lee, Stan, writer, and Steve Ditko, artist. \textit{Amazing Fantasy} no. 15. Marvel Comics, 1962.
  \item Lapham, David, writer, and Tony Harris, artist. \textit{Spider-Man: With Great Power\ldots} no. 1--5. Marvel Comics, 2008.
  \item Slott, Dan, writer, and Stefano Caselli, artist. \textit{Avengers: The Initiative} no. 1--6. Marvel Comics, 2007.
\end{dossieritems}
}

\dossierbox{\observationstitle}{%
This is a fictional classroom role-play; the Avengers Academy, characters, incident, and mission details are invented. Keep the activity discussion-based: do not use hazardous props or ask recruits to reenact an emergency. Adapt roles and materials to students' needs, and reinforce that real emergencies require following local safety procedures and seeking qualified help.
}
